\documentclass{article}

\usepackage{amsmath}
\usepackage{amssymb}
\usepackage{amsthm}
\usepackage{stmaryrd}
\usepackage{ mathrsfs }
\usepackage{ mathtools}
\usepackage{enumerate}
\usepackage{bbm}
\usepackage{lipsum}
\usepackage{fancyhdr}
\usepackage{stix}
\usepackage{tikz-cd} 
\usepackage{qworld}
\usetikzlibrary{tqft}
\tikzset{commutative diagrams/.cd,
mysymbol/.style={start anchor=center,end anchor=center,draw=none}
}
\newcommand\comsymb[2][\circlearrowleft]{%
  \arrow[mysymbol]{#2}[description]{#1}}

\newtheorem{theorem}{Theorem} [section] 
\newtheorem{proposition}{Proposition}[section] 
\newtheorem{definition}{Definition}[section] 
\newtheorem{lemma}{Lemma}[section] 
\newtheorem{notation}{Notation}[section] 
\newtheorem{remark}{Remark}[section] 
\newtheorem{corollary}{Corollary} [section] 
\newtheorem{terminology}{Terminology}[section] 
\newtheorem{fact}{Fact}[section] 
\newtheorem{claim}{Claim}[proposition]
\newtheorem{example}{Example}[section] 
\newtheorem{exercise}{Exercise}[section] 
\numberwithin{equation}{theorem}

\DeclareMathOperator{\Hilb}{Hilb}
\DeclareMathOperator{\Frob}{Frob}
\DeclareMathOperator{\FA}{FA}
\DeclareMathOperator{\Cob}{Cob}
\DeclareMathOperator{\rev}{rev}
\DeclareMathOperator{\TQFT}{TQFT}
\DeclareMathOperator{\st}{s.t.}
\DeclareMathOperator{\rad}{rad}
\DeclareMathOperator{\chr}{char}
\DeclareMathOperator{\alggrp}{alg-grp}
\DeclareMathOperator{\Simple}{Simple}
\DeclareMathOperator{\ab}{ab}
\DeclareMathOperator{\len}{len}
\DeclareMathOperator{\coker}{coker}
\DeclareMathOperator{\triv}{triv}
\DeclareMathOperator{\ev}{ev}
\DeclareMathOperator{\tr}{tr}
\DeclareMathOperator{\ad}{ad}
\DeclareMathOperator{\ord}{ord}
\DeclareMathOperator{\Aut}{Aut}
\DeclareMathOperator{\Vect}{Vect}
\DeclareMathOperator{\End}{End}
\DeclareMathOperator{\Irr}{Irr}
\DeclareMathOperator{\Map}{Map}
\DeclareMathOperator{\Mor}{Mor}
\DeclareMathOperator{\Stab}{Stab}
\DeclareMathOperator{\Ind}{Ind}
\DeclareMathOperator{\GL}{GL}
\DeclareMathOperator{\SL}{SL}
\DeclareMathOperator{\SO}{SO}
\DeclareMathOperator{\Orth}{O}
\DeclareMathOperator{\Sp}{Sp}
\DeclareMathOperator{\Spin}{Spin}
\DeclareMathOperator{\Ad}{Ad}
\DeclareMathOperator{\Rep}{Rep}
\DeclareMathOperator{\trdeg}{trdeg}
\DeclareMathOperator{\Frac}{Frac}
\DeclareMathOperator{\rank}{rank}
\DeclareMathOperator{\Hom}{Hom}
\DeclareMathOperator{\Proj}{Proj}
\DeclareMathOperator{\Homk}{Hom_{k-alg}}
\DeclareMathOperator{\Ker}{ker}
\DeclareMathOperator{\Image}{im}
\DeclareMathOperator{\image}{im}
\DeclareMathOperator{\Coker}{Coker}
\DeclareMathOperator{\Dom}{Dom}
\DeclareMathOperator{\grad}{grad}
\DeclareMathOperator{\Span}{Span}
\DeclareMathOperator{\Spec}{Spec}
\DeclareMathOperator{\MaxSpec}{MaxSpec}
\DeclareMathOperator{\interior}{int}
\DeclareMathOperator{\supp}{supp}
\DeclareMathOperator{\id}{id} 
\DeclareMathOperator{\op}{op} 
\DeclareMathOperator{\sgn}{sgn}
\DeclareMathOperator{\Mat}{Mat}
\DeclareMathOperator{\Adj}{Adj}
\DeclareMathOperator{\ConcAffVar}{ConcAffVar}
\DeclareMathOperator{\AbsAffVar}{AbsAffVar}
\DeclareMathOperator{\pr}{pr}
\DeclareMathOperator{\Mod}{Mod}
\DeclareMathOperator{\Comod}{Comod}
\DeclareMathOperator{\Free}{Free}
\DeclareMathOperator{\indec}{indec}
\DeclareMathOperator{\Res}{Res}
\DeclareMathOperator{\Set}{Set}
\DeclareMathOperator{\SVect}{SVect}
\DeclareMathOperator{\FSet}{FSet}
\DeclareMathOperator{\Top}{Top}
\DeclareMathOperator{\Sk}{Sk}
\DeclareMathOperator{\coev}{coev}
\DeclareMathOperator{\Ver}{Ver}
\DeclareMathOperator{\ComAlg}{ComAlg}
\DeclareMathOperator{\Grp}{Grp}
\DeclareMathOperator{\PGL}{PGL}
%\newcommand*{\name}[\num_arguments][default values]{{\color{#1}\Large #2}}
\newcommand{\defeq}{\vcentcolon=}
\newcommand{\norm}[1]{\Vert #1 \Vert}
\newcommand{\opNorm}[2]{\norm{#1}_{#2\to#2}}
\newcommand{\normL}[3]{\norm{#1}_{L^{#2}(#3)}}
\newcommand{\N}[0]{\mathbb{N}}
\newcommand{\m}[0]{\mathfrak{m}}
\newcommand{\R}[0]{\mathbb{R}}
\newcommand{\Ouv}[0]{\mathcal{O}}
\newcommand{\Z}[0]{\mathbb{Z}}
\newcommand{\C}[0]{\mathbb{C}}
\newcommand{\bQ}[0]{\mathbb{Q}}
\newcommand{\F}[0]{\mathbb{F}}
\newcommand{\G}[0]{\mathbb{G}}
\newcommand{\A}[0]{\mathbb{A}}
\newcommand{\Para}[0]{\mathbb{P}}
\newcommand{\K}[0]{\mathbb{K}}
\newcommand{\cat}[0]{\mathscr{C}}
\newcommand{\catd}[0]{\mathscr{D}}
\newcommand{\M}[0]{\mathbb{M}}
\newcommand{\Lie}[0]{\mathcal{L}}
\newcommand{\torus}[0]{\mathbb{T}}
\newcommand{\g}[0]{\mathfrak{g}}
\newcommand{\tl}[0]{\mathfrak{t}}
\newcommand{\gl}[0]{\mathfrak{gl}}
\newcommand{\spl}[0]{\mathfrak{sl}}
\newcommand{\so}[0]{\mathfrak{so}}
\newcommand{\spg}[0]{\mathfrak{sp}}
\newcommand{\sep}[0]{\:|\:}
\newcommand{\card}[1]{\vert #1\vert}
\newcommand{\smat}[1][n]{\Mat_{#1\times #1}}
\newcommand{\I}[1][n]{\mathcal{I}}
\newcommand{\disum}[1][n]{\stackrel{\subseteq}{\oplus}}
\newcommand{\summand}[0]{\stackrel{\subseteq}{\oplus}}
\newcommand{\fib}[1]{%
  \mathbin{\mathop{\times}\limits_{#1}}%
}
\newcommand{\tens}[1]{%
  \mathbin{\mathop{\otimes}\displaylimits_{#1}}%
}

\title{Representation Theory 2}
\author{So Murata}
\date{WiSe 25/26, University of Bonn}


\begin{document}
\maketitle
% \subsection{Negligible Morphisms}
% \begin{definition}
%     Let $a$ be a pivotal structure for the category $\cat$. $a$ is said to be spherical if for any morphism $f$ of $\cat$,
%     \begin{equation*}
%         \tr^L(f) = \tr^R(f)=0.
%     \end{equation*}
% \end{definition}

% \begin{example}
% $\Vect,\Rep G,\Z/2\Z-\Vect$ admit spherical structures.
% \end{example}

% \begin{definition}
%     \label{definition_negligible_morphisms}
%     Let $\cat$ be a rigid monoidal category with a spherical structure $a$. A morphism $f:X\to Y$ in $\cat$ is negligible if for any $g\in\cat(Y,X)$, we have $\tr(gf) = 0$. 
%     \par The collection of all negligible morphisms is denoted by $\mathcal{N}$.
% \end{definition}
% \begin{remark}
%     By the swapping property of $\tr$, we can rephrase Definition \ref{definition_negligible_morphisms} can be restated by $\tr(fg) = 0$ for any $g:Y\to X$.
% \end{remark}
% \begin{lemma}
%     $\mathcal{N}$ is a tensor ideal.
% \end{lemma}
% \begin{proof}
%     For $f\in\mathcal{N},f',g$ be morphisms then,
%     \begin{equation*}
%         \tr(gf'f)=0,\tr(gff') = \tr(f'gf)=0,
%     \end{equation*}
%     whenever these expressions make sense. Thus $f'f,ff'\in\mathcal{N}$. Consider,
%     \begin{equation*}
%         \tr(g(f\otimes f'))=\q{\qcap[p=4]\n\qbox[name={$f$}]\qbox[name={$f'$}]\qwire\qwire\n\qbox[name={$f$},p=2]\qwire\qwire\n\qbox[name={$a$}]\qbox[name={$a$}]\qwire\qwire\n\qcup[p=4]\qbb[4,10]\qbb[6]}=\q{\qcap\n\qbox[name={$f$}]\qwire\n\qbox[name={$g'$}]\qwire\n\qbox[name={$a$}]\qwire\n\qcup}=\tr(g'f) = 0,
%     \end{equation*}
%     where $g'$ is the morphism corresponding to the diagram bounded by dotted lines. Thus $f\otimes f'\in\mathcal{N}$. Similarly,
%     \begin{equation*}
%         \tr(g(f'\otimes f))=\q{\qcap[p=4]\n\qwire\qwire\qbox[name={$a^{-1}$}]\qbox[name={$a^{-1}$}]\n\qwire\qwire\qbox[name={$f'$}]\qbox[name={$f$}]\n\qwire\qwire\qbox[name={$g$},p=2]\n\qcup[p=4]\qbb[4,11]\qbb[12]} =\tr(g'f) = 0,
%     \end{equation*}
%     where $g'$ is the morphism corresponding to the diagram bounded by dotted lines. We now have $f'\otimes f\in\mathcal{N}$. For the last equality, we used $a$ is spherical.
% \end{proof}
% \begin{lemma}
%     \label{lemma_decomposition_endomorphism}
%     Let $\cat$ be an additive idempotent complete $\K$-linear Hom-finite category where $\K$ is algebraically closed. Then for any $X\in \indec(\cat)$ and $f\in\End_{\cat}(X)$, there is $\lambda\in\K$ and a nilpotent $n\in\End_\cat(X)$ such that $f=\lambda+n$. 
%     \par For any $X,Y\in\indec(\cat)$ and $f\in\cat(X,Y)$ $f$ is not isomorphism if and only if for any $g\in\cat(Y,X)$, $gf$ is a nilpotent.
% \end{lemma}
% \begin{proof}
%     By Proposition \ref{proposition_krull_schmidt}, $R\defeq\End_\cat(X)$ is a local ring and as $\K=\overline{\K}$ we have,
%     \begin{equation*}
%         R/\rad(R)=\K.
%     \end{equation*}
%     Since $\rad(R)$ consists of exactly the nilpotent elements, we have the statement.
%     \par For the second assertion, suppose $f$ is an isomorphism then $f{-1}f$ is not a nilpotent. Therefore, it is enough to show that $f$ is not an isomorphism then $gf$ is a nilpotent. Suppose $gf$ is not a nilpotent, then by the previous part it is invertible. Indeed, we can take,
%     \begin{equation*}
%         \sum_{i=0}^\infty (-1)^in^i.
%     \end{equation*}
%     As $n$ is a nilpotent, this sum is finite and defines an inverse of $1+n$. Therefore $gf$ is invertible. Set $g'\defeq g(fg)^{-1}$, then $fg'=\id$.
%     Let $e\defeq g'f$. Then $e$ is idempotent. In particular, it is an idempotent element of an indecomposable object. As $\End_\cat(X)$ is local by Proposition \ref{proposition_krull_schmidt} and idempotent elements of a local ring are trivial. We have,
%     \begin{equation*}
%         e=
%         \begin{cases}
%             0\Rightarrow f=fg'f = fe\Rightarrow (fg) = 0\Rightarrow X\cong Y\cong 0,\\
%             1\Rightarrow \text{$f$ is an isomorphism}.
%         \end{cases}
%     \end{equation*}
%     This completes the proof.
% \end{proof}
% \begin{lemma}
%     \label{lemma_negligible_equivalence}
%     In the situation of Lemma \ref{lemma_decomposition_endomorphism}, assume $(\cat,a)$ is a spherical structure. Consider $X,Y\in\indec(\cat)$ and $f:X\to Y$. Then $f$ is negligible if and only if either it is not an isomorphism or $\dim X=0=\dim Y$.
% \end{lemma}
% \begin{proof}
%     For any $g:Y\to X$, there exists $\lambda_g\in\K,n_g\in\End_\cat(Y)$ where $n_g$ is a nilpotent. such that 
%     \begin{equation*}
%         fg = \lambda_g+n_g.
%     \end{equation*}
%     Then 
%     \begin{equation*}
%         f\in\mathcal{N}\Leftrightarrow \forall g:Y\to X,\tr(fg)=0\Leftrightarrow \lambda_g\tr(\id_Y)=0\Leftrightarrow\lambda_g=0\lor\dim Y=0.
%     \end{equation*}
%     If $\lambda_g=0$, then $fg$ is a nilpotent and by Lemma \ref{lemma_decomposition_endomorphism}, $f$ is not an isomorphism. If $\lambda_g\not=0$ and $\dim Y=0$, then $f$ is an isomorphism thus $\dim X=0$.
% \end{proof}
% \begin{theorem}
%     In the situation of Lemma \ref{lemma_decomposition_endomorphism}, $\cat/\mathcal{N}$ is semisimple with simple objects given by 
%     \begin{equation*}
%     \{X\sep X\in\indec(\cat),\dim(X)\not=0\}.
%     \end{equation*}
% \end{theorem}
% \begin{proof}
%     For any indecomposable $X$ with $\dim X\not=0$. Decomposition in Lemma \ref{lemma_decomposition_endomorphism} and Lemma \ref{lemma_negligible_equivalence}  tell us that 
%     \begin{equation*}
%         \cat(X,Y)/\mathcal{N}(X,Y) \cong\begin{cases}
%             \K,X\cong Y,\\
%             0,X\not\cong Y.
%         \end{cases}
%     \end{equation*}
%     This shows that $X$ is simple. 
% \end{proof}
% \begin{remark}
%     If $\cat$ is semisimple then $\mathcal{N}=\{0\}$. Indeed if $f\in\mathcal{N}(X,Y)$, then direct sums of $X,Y$ gives a morphism $f:X_i\to Y_j$ where $X_i,Y_j$ are simple objects which are summands of $X,Y$, respectively. $f$ is not isomorphism as by assumption thus it is $0$.
% \end{remark}
% \begin{remark}
%     $\cat/\mathcal{N}$ is calledthe maximal semisimple quotient. This has an induced spherical structure and the functor $\cat\to\cat/\mathcal{N}$ preserves the dimensions.
% \end{remark}
% \begin{example}
% Suppose $\K$ is algebraically closed and $\chr\K=p$, we will define.
% \begin{equation*}
%     \Ver_p\defeq\Rep_\K(\Z/p\Z)/\N.
% \end{equation*}
% This category has indecomposable objects $\overline{X}_1,\cdots,\overline{X}_{p-1}$ where for each $i=1,\cdots,p-1$, $X_i$ corresponds to a Jordan block of size $i$.
% \par For $p=2$, indecomposable objects of $\Ver_2$ is $X_1=\K$. Therefore, $\Ver_2\cong\Vect_\K$ as linear monoidal categories. 
% \par For $p=3$, indecomposable objects are $\{X_1=\K,X_2\}$. From Example \ref{example_indecomposable_jordan_block}, $X_2\otimes X_2 = X_1\oplus X_3$ where we conventionally write $X_3=0$. We note that $\Ver_3\cong\SVect$. 
% \end{example}
% \begin{example}
%     Let us examine another important example of $\Ver_p$ in the view of Example \ref{example_indecomposable_jordan_block}. Namely, the case where $p=5$. 
%     Let us denote the indecomposables by $\{X_1=1,Y_2,\cdots,Y_4\}$. Then 
%     \begin{center}
% \begin{tabular}{|l|l|l|l|}
% \hline
%       & $Y_2$     & $Y_3$     & $Y_4$ \\ \hline
% $Y_2$ & $Y_1+Y_3$ & $Y_2+Y_4$ & $Y_3$ \\ \hline
% $Y_3$ & $Y_2+Y_4$ & $Y_1+Y_3$ & $Y_2$ \\ \hline
% $Y_4$ & $Y_3$     & $Y_2$     & $Y_1$ \\ \hline
% \end{tabular}
%     \end{center}
%     wwhere we calculate rows besides the first one as follows.
%     \begin{align*}
%         Y_3^2  & = Y_3(Y^2_2-1)=(Y_2+Y_4)Y_2-Y_3,\\
%         &= Y_1+2Y_3-Y_3 = Y_1+Y_3.\\
%         Y_3Y_4 & = (Y_2^2-1)Y_4 = Y_2+Y_4-Y_4 = Y_2.\\
%         Y_4^2 & = Y_4(Y_2Y_3-Y_2) = Y_3^2-Y_3^2 = Y_1.
%     \end{align*}
%     \label{example_ver_5}
% \end{example}
% \begin{remark}
%     In general we have $X_{p-1}X_i=X_{p-i}$. Thus $\langle 1,X_{p-1}\rangle_{\oplus}$ is closed under $\otimes$ and is isomorphic to $\SVect$. Furthermore $\langle 1,X_3,X_5,\cdots\rangle_{\oplus}$ is closed under $\otimes$ as well. 
%     We define,
%     \begin{equation*}
%         \Ver_p^+\defeq \langle 1,X_3,X_5,\cdots\rangle_{\oplus}.
%     \end{equation*}
%     We can check in Example \ref{example_ver_5} that $\Ver_5=\langle 1,X_3\rangle_{\oplus}$ is indeed closed under $\otimes$.
%     \par Suppose we have a functor $F:\Ver_5^+\to\SVect$ which is linear monoidal. Let $Y\defeq F(X_3)$. Then $Y\otimes Y=Y\oplus 1$. This tells us that for $d=\dim_\K Y$,
%     \begin{equation*}
%         d^2=d+1.
%     \end{equation*}
%     As $d$ is integer, we conclude there is not reasonable such functor $F$.
% \end{remark}
%last lecture below
% \begin{theorem}
%   For any tuple $(\alpha,u)$, we have $\Cob_{\alpha,u}$ is Hom-finite and $\End(1)\cong \K$. Conversely if $\cat$ is Hom-finite and $\End(1)\cong \K$, then for any symmetric monoidal functor $F:\Cob\to\cat$, there is a Frobenius algebra datum $(\alpha,u)$ such that 
%   \begin{center}
%     \begin{tikzcd}
% \Cob \arrow[rr, "F"] \arrow[rd, two heads] &                                                  & \cat \\
%                                            & {\Cob_{\alpha,u}} \arrow[ru, "\exists"', dotted] &     
% \end{tikzcd}
%   \end{center} 
% \end{theorem}

% \begin{proof}
%   In the first statement, to show Hom-finiteness, we show that all Hom-spaces are spanned by cobordisms of genus less than $d-1$ where $d$ is the degree of $u$. as we have the relation,
%   \begin{equation*}
%     \text{Insert cobordisms diagrams here}
%   \end{equation*}
%   For $\End(1)=\K$, it is spanned by cobordisms with a boundary 
%   \begin{equation*}
%     \text{Insert cobordisms diagrams here} = \alpha_i\in\K
%   \end{equation*}
%   To show the second statement, notice that $F([1])\in\Frob(\cat)$. Take 
%   \begin{equation*}
%     \phi\defeq F\left(\text{cobordism diagram}\right).
%   \end{equation*}
%   As $\cat$ is Hom-finite, there is $u\in\K[x],u(\phi)=0$. As $\End_\cat(1)\cong \K$, there is $\alpha_i\in\K$ such that 
%   \begin{equation*}
%     F\left(\text{cobordism diagram}\right) = \alpha_i.
%   \end{equation*}
%   This mesn $F$ factors via $\Cob_{\alpha,u}$.
% \end{proof}
% \begin{example}
%     For $\cat=\Rep G$ where $G$ is finite and $X$ be a $G$-set. We have,
%     \begin{equation*}
%         F:\Cob\to\Rep G,[1]\mapsto\K.
%     \end{equation*}
%     In this case,
%     \begin{align*}
%         \phi=F\left(\text{cobordisms diagra}\right) = \q{\qmul{$\K X$}\n\qcomul{$\K X$}} ,\\
%         & = (X\stackrel{\Delta}{\mapsto}X\otimes X\stackrel{\mu}{\mapsto}X) = \id, x\in X.
%     \end{align*}
%     So $u(X) = X-1$, we have,
%     \begin{align*}
%         \alpha_i & = F\left(\text{cobordism diagram}\right),\\
%         & = F\left(\text{cobordism diagram}\right)F\left(\text{cobordism diagram}\right)F\left(\text{cobordism diagram}\right),\\
%         & = \varepsilon\phi^i,\\
%         \eta & \varepsilon\eta = (1\stackrel{\eta}{\mapsto}\sum_{x\in X}x\stackrel{\varepsilon}{\mapsto}X).
%     \end{align*}
%     So $\Cob\to\Rep G$ factors via 
%     \begin{equation*}
%         \Cob_{\underbrace{(\vert X\vert,\vert X\vert,\cdots)}_{=\alpha},\underbrace{x-1}_{u}}.
%     \end{equation*}
%     \begin{tikzpicture}
%     \pic[tqft/cylinder, at={(1,0)}, fill=green, genus=1];
% \end{tikzpicture}
% \end{example}
% \begin{definition}
%     $\forall t\in\K$, we define,\begin{equation*}
%         \underline{\Rep}S_t\defeq \Cob_{((t,t,\cdots),X-1)}
%     \end{equation*}
% \end{definition}

% \begin{proposition}
%     Assume $\chr(\K)=0$ and $n\in\Z_{\geq0}$. Then there is a $\K$-linear, $\otimes$ symmetric functor 
%     \begin{equation*}
%         \underline{\Rep}S_n\to\Rep S_n,
%     \end{equation*}
%     such that 
%     \begin{equation*}
%         [1]\mapsto V=\K^n=\K\{1,\cdots,n\},
%     \end{equation*}
%     which is full, essentially surjective and the semisimplification factor.
% \end{proposition}
% \begin{proof}
%     For fullness, consider 
%     \begin{align*}
%         \Hom_{S_n}(V^{\otimes m},V^{\otimes n})&\cong \Hom_\K(\K X^m,\K X^n)^{S_n},\\
%         & \cong \Hom_\K(\K,\K X^{m+n})^{S_n},\\
%         & \cong (\K X^{m+n})^{S_n},\\
%         & = \K\left\{\sum S_ny\sep y\in\K X^{m+n}\right\},\\
%         & = \K\left\{\sum\right\}_{y\in \K X^{m+n}\sep\text{such that $y_i=y_j$ where $i,j$ (missing) in the same com of $P$}},
%         &\cong\K\{\text{set partition of $\{1,\cdots,m+m'\}$}\}
%     \end{align*}
%     For essentially surjectivenes, as $S_n$ is a finite group, any irreducible representation $W$ of $S_n$ appears as a direct summand in $V^{\otimes m}$ (there exists such $m$) as $V$ is faithful then $W=(V^{\otimes m},e)$ where $e$ is an idempotent. As $F$ is full, $e$ is with image. $W$ is in the image.
%     \par For the semisimplification, $\Rep S_n$ is semisimple, there is no negligible morphisms in $\Rep S_n$. This means the negligible morphisms in $\underline{\Rep}S_t$ must be contained in the kernel of $F$. Thus,
%     \begin{center}
%         \begin{tikzcd}
% \underline{\Rep}S_t \arrow[rr] \arrow[rd] &          & \underline{\Rep}S_t/\mathcal{N} \arrow[ld, "\exists", dashed] \\
%                                           & \Rep S_t &                                                              
% \end{tikzcd}
%     \end{center}
%     Assume $\mathcal{N}\subsetneq\ker F$ then $\G$ is a proper quotinet. Thus $\underline{\Rep}_G/\mathcal{N}$ has a proper tensor ideal. This contradicts to the semisimplicity of the tensor category.
% \end{proof}
% \begin{remark}
%     This is why we call $\Rep S_t$ interpolation category. For $t\not\in\Z_{\geq0}$ then $\underline{\Rep}S_t$ is semisimple and hence, a tensor category not of moderate growth.
% \end{remark}
% \subsection{Abelian realization}
% Recall $\Cob_{\alpha,u}$ has a $(\otimes,\text{Sym},\K)$-functor to a tensor category only if those that factor via $\Cob_{\alpha,u}$ for some Frobenius algebra datum $(\alpha,u)$. In particular,
% \begin{equation*}
%     \forall j,\sum_{i\geq0}u_i\alpha_{i+j}=0.
% \end{equation*}
% \begin{definition}
%     For all $\alpha=(\alpha_i)$ in $\K$, $u_\alpha\defeq$ monic polynomial of minimal degree such that $(\alpha,u)$ is a Frobenius datum that is 
%     \begin{equation*}
%         \forall j\geq0, \sum_{i\geq0}u_i\alpha_{i+j}=0.
%     \end{equation*}
%     We then define,
%     \begin{equation*}
%         \defeq\Cob_\alpha \defeq\Cob_{\alpha,u_\alpha}, \FA_{\ab} = \{\alpha\sep(\alpha,X^2(X_{r_1})\cdots(X_{r_k}))\}.
%     \end{equation*}
%     is a Frobenius algebra datum for distinct $r_1,\cdots,r_k\in\K$.
% \end{definition}
% In this subjection, we will assume $\K=\overline{\K}$ and $\chr\K=0$. 
% \begin{theorem}
%     For all $\alpha$ there is $\Cob_\alpha\to\cat$ where $\cat$ is a tensor category, if and only if $\alpha\in \FA_{\ab}$.
% \end{theorem}
% \begin{proof}$\:$
%     \par\textbf{$\Rightarrow$} As before we take $\phi=F(\text{some cobordism diagram})$. As $\cat$ is a Hom-finite, there is $\hat{u}$ such that $\hat{u}(\phi)=0$. Set,
%     \begin{equation*}
%         u(x) = \prod_{\text{missing}}(x-r).
%     \end{equation*}
%     Then $u(\phi)$ is nilpotent. As $\cat$ is a tensor category, we have,
%     \begin{equation*}
%         0=\sum_{i\geq0}u_i,\underline{\tr(\phi)}_{\text{missing}} = \sum_{i\geq0}()_{i,\alpha_i}.
%     \end{equation*}
% \end{proof}
% \begin{proposition}
%     For all $\alpha\in\FA_{\ab}$, there exist $t_1,\cdots,t_k,\alpha_0,\alpha_1\in\K$ such that 
%     \begin{equation*}
%         \Cob_\alpha\cong \Cob({\alpha_0,\alpha_1,0,0,\cdots})\boxtimes\underline{\Rep}S_{t_1}\boxtimes\cdots\boxtimes\underline{\Rep}S_{t_k}.
%     \end{equation*}
%     where $\boxtimes$ is the tensor product of categories. Explicitly, 
%     \begin{equation*}
%         \cat\boxtimes\catd \defeq(\cat\fib{\K}\catd)^{\oplus,\summand}.
%     \end{equation*}
%     $\cat\boxtimes\catd$ has objects $X\boxtimes Y$ where $X\in\cat,Y\in\catd$ such that 
%     \begin{equation*}
%         \Hom(X\boxtimes Y,X'\boxtimes Y')\cong\Hom_\cat(X,X')\tens{\K}\Hom_\catd(Y,Y').
%     \end{equation*}
% \end{proposition}
% \begin{proof}
%     The idea is, each functor on the right hand side has a canonical commutative Frobenius algebra $A_i$. Set,
%     \begin{equation*}
%         A\defeq A_1\boxtimes1\boxtimes\cdots\boxtimes 1\oplus\cdots\oplus1\boxtimes\cdots\boxtimes 1\boxtimes A_{k+1}
%     \end{equation*}
%     is a commutatie Frobenius algebra. This $A$ is the correct Frobenius algebra datum $(\alpha,u_\alpha)$ so there is $\Cob_\alpha\to$ the right hand side.
%     \par Then show this is an equivalence.
% \end{proof}
% \begin{proposition}
%     There is a $\K$-linear symmetric tensor category functor $\Cob_{(\alpha_0,\alpha_1,0,\cdots,)}\to\cat$ where $\cat$ is a tensor category. For example,
%     \begin{enumerate}[1).]
%         \item $\alpha_0=\alpha_1=0$ then $\cat=\Vect$.
%         \item $\alpha_1=0\not=\alpha_0$, then $\cat=\Rep\Grp(1|2)$. 
%         \item $\alpha_1\not=0$, then 
%         \begin{equation*}
%             \cat = \begin{cases}
%                 \underline{\Rep}O_t,\quad \alpha_1\not\in\Z_{\geq0},\\
%                 \Rep O_n,\quad \alpha=n\in\Z_{\geq0},\\
%                 \Rep \underbrace{O\Sp(m|2)}_{\text{\makebox[0pt][c]{Orthosymplectic supergroup}}},\quad \alpha_1=m-2n
%             \end{cases}
%         \end{equation*}
%         The first one is an analogue of oriented Brauer category $\cat_t$ without the orientation.
%     \end{enumerate}
% \end{proposition}
% \begin{corollary}
%     $\Leftarrow$ of the theorem holds (mention).
% \end{corollary}
% \begin{remark}
%     This provides a universal abelian version of the theory of TQFTs (functors from cobordisms).
% \end{remark}
% \begin{example}
%     $\alpha_1=\alpha_0=0$, then $u_\alpha=1$. Then in $\Cob_\alpha$, any non-empty cobordism factors through cylinder. And cylinder is $0$. Thus any cobordism with a boundary circle is $0$. $\alpha_i=0$ for all $i$, thus all non-empty cobordism is $0$.
%     Thus $\Cob_\alpha=\Vect$ where $[0]\leftrightarrow\K,(\id_{[0]}=\empty\not=0)$. 
% \end{example}
% \begin{example}
%     $\alpha_1=0$, then $u_\alpha=X$ in $\Cob_\alpha$, $\text{cobordism diagram}=0$. Thus any cobordism of genus more than $1$ is isomorphic.
% \end{example}
% So we can represent morphisms by partition diagrams. Also any cobordism with at least one incoming and one outgoing boundary cycle is in $\mathcal{N}$.
% \begin{example}
%     For all $m,n\geq0$, consider,
%     \begin{equation*}
%         \gl(m|n) = \left\{\begin{pmatrix}
%             A & B\\ C & D
%         \end{pmatrix}
%         \right\} = 1.
%     \end{equation*}
%     where $A$ is $m\times m$, $B$ is $m\times n$, $C$ is $n\times m$, and $D$ is $n\times n$. and,
%     \begin{equation*}
%        \gl(m|n)_D = \left\{\begin{pmatrix}
%             A & \:\\ \: & D
%         \end{pmatrix}
%         \right\} = 1,\gl(m|n)_J = \left\{\begin{pmatrix}
%             \: & B\\ C & \:
%         \end{pmatrix}
%         \right\} = 1.
%     \end{equation*}
%     $\gl(m|n)$ is a Lie super algebra with the (missing),
%     \begin{equation*}
%         [A,B] = AB-(-1)^{(A|B)}BA,
%     \end{equation*}
%     for all $A,B$ homogeneous in $\gl(m|n)$. 
%     \par Consider,
%     \begin{equation*}
%         \g(1|2)\defeq\left\{\begin{pmatrix}
%             0 & a & b \\
%             b & \: & \:\\
%             -a & \: & A
%         \end{pmatrix}
%         \:\Bigg{|}\:a,b\in\C,A\in \spl_2(\C)\right\}\subseteq\gl(1|2).
%     \end{equation*}
%     For $\Rep\mathfrak{osp}(1|2)\ni V=\K^{1|2}$ by matrix vector-multiplication.
%     \par $\Rep\mathfrak{osp}(1|2)^+=\langle V\rangle_{\otimes,\oplus,\summand}\subseteq\Rep\mathfrak{osp}(1|2)$.
% \end{example}
% \begin{remark}
%     $\Rep\mathfrak{osp}(1|2)^+$ is semisimple and any representation is determined by $X=X_{\chi_2}$. 
%     \par For example $X(V) = q+1+q^{-1}$ as 
%     \begin{equation*}
%         \begin{pmatrix}
%             0 & \: & \:\\
%             \: & 1 & \:\\
%             \: & \: & -1
%         \end{pmatrix}
%         \oplus\begin{pmatrix}
%             x_0\\ x_+\\ x_-.
%         \end{pmatrix}
%     \end{equation*}
%     \begin{align*}
%         X(V^{\otimes 2}) &= (q+1+q^{-1}),\\
%         & = q^2+2q+3+2q^{-1}+q^{-2},\\
%         & = \underbrace{(1)}_{X(1)}+\underbrace{(q+f+q^{-1})}_{X(V)}+(q^{-2}+q^+f+q^{-1}+q^{})
%     \end{align*}
% \end{remark}
\subsection{Braidings}
\begin{definition}
  A braiding/braided category is a pair $(\cat,c)$ where $\cat$ is a monoidal category and $c = (c_{X,Y}:X\otimes Y\to Y\otimes X)$ is a ntural isomorphism such that for any $X,Y,Z\in\cat$, we have the commudatie diagrams.
  \begin{center}
\begin{tikzcd}
X\otimes Y\otimes Z \arrow[rr, "{c_{X,Y\otimes Z}}"] \arrow[rd, "{c_{X,Y}\otimes\id_Z}"']  &                                                           & Y\otimes Z\otimes X \\
                                                                                           & Y\otimes X\otimes Z \arrow[ru, "{\id_Y\otimes c_{X,Z}}"'] &                     \\
X\otimes Y\otimes Z \arrow[rr, "{c_{X\otimes Y,Z}}"] \arrow[rd, "{\id_X\otimes c_{Y,Z}}"'] &                                                           & Z\otimes X\otimes Y \\
                                                                                           & X\otimes Z\otimes Y \arrow[ru, "{c_{X,Z}\otimes\id_Y}"']  &                    
\end{tikzcd}
  \end{center}
  The braiding $c$ is called symmetric if $c^2=\id$. That is for any $X,Y\in\cat$ we hae,
  \begin{equation*}
    c_{X,Y} = c_{Y,X}^{-1}.
  \end{equation*}
\end{definition}
\begin{lemma}
    Let $(\cat,c)$ be a braided category. We have the following statements.
    \begin{enumerate}[1).]
        \item For $X\in\cat$, 
        \begin{equation*}
            c_{X,1}=\lambda_X^{-1}\rho_X,c_{1,X}=\rho_{X}^{-1}\lambda_X = c_{X,1}^{-1}.
        \end{equation*}
        \item Young-Baxter equation holds. That is,\begin{equation*}
            (c\otimes\id)(\id\otimes c)(c\otimes\id) = (\id\otimes c)(c\otimes\id)(\id\otimes c).
        \end{equation*}
        This is represented in the commutative diagram as follows.
        \begin{center}
            \begin{tikzcd}
Y\otimes X\otimes Z \arrow[rr, "\id\otimes c"]                          &  & Y\otimes Z\otimes X \arrow[d, "c\otimes\id"]   \\
X\otimes Y\otimes Z \arrow[u, "c\otimes\id"] \arrow[d, "\id\otimes c"'] &  & Z\otimes Y\otimes X                            \\
X\otimes Z\otimes Y \arrow[rr, "c\otimes \id"']                         &  & Z\otimes X\otimes Y \arrow[u, "\id\otimes c"']
\end{tikzcd}
        \end{center}
    \end{enumerate}
\end{lemma}
\begin{proof}
    In the view of Theorem \ref{theorem_maclane_coherence}, and use $c$ is a natural isomorphism, one has,
    \begin{center}
        \begin{tikzcd}
X\otimes(1\otimes 1) \arrow[rrrr, "\alpha^{-1}"] \arrow[rd, "\id_X\otimes\lambda" description] \arrow[dddd] &                                                  &                                                                               &   & (X\otimes 1)\otimes 1 \arrow[llld, "\rho" description] \arrow[d, "c\otimes\id_1"] \\
                                                                                                            & X\otimes 1 \arrow[dd] \arrow[r, "c" description] & 1\otimes X \arrow[dd, "c^{-1}" description] \arrow[rd, "\lambda" description] &   & (1\otimes X)\otimes 1 \arrow[dd, "\alpha"] \arrow[ll, "\rho" description]         \\
                                                                                                            &                                                  &                                                                               & X &                                                                                   \\
                                                                                                            & 1\otimes X                                       & X\otimes 1 \arrow[l, "c" description] \arrow[ru, "\rho" description]          &   & 1\otimes(X\otimes1) \arrow[d, "\id_1\otimes c"] \arrow[ll, "\lambda" description] \\
(1\otimes 1)\otimes X \arrow[ru, "\lambda\otimes\id_X" description] \arrow[rrrr, "\alpha"']                 &                                                  &                                                                               &   & 1\otimes(1\otimes X) \arrow[lllu, "\lambda" description]                         
\end{tikzcd}
    \end{center}
    This proves the first assertion. For the second assertion, 
    \begin{center}
        \begin{tikzcd}
Y\otimes X\otimes Z \arrow[rr, "{\id_Y\otimes c_{X,Z}}"]                                                                                  &  & Y\otimes Z\otimes X \arrow[d, "{c_{Y,Z}\otimes\id_X}"]   \\
X\otimes Y\otimes Z \arrow[rru, "{c_{X,Y\otimes Z}}" description] \arrow[u, "{c_{X,Y}\otimes\id_Z}"] \arrow[d, "{\id_X\otimes c_{Y,Z}}"'] &  & Z\otimes Y\otimes X                                      \\
X\otimes Z\otimes Y \arrow[rr, "{c_{X,Y}\otimes \id_Y}"'] \arrow[rru, "{c_{X,Z\otimes Y}}" description]                                   &  & Z\otimes X\otimes Y \arrow[u, "{\id_Z\otimes c_{X,Y}}"']
\end{tikzcd}
    \end{center}
\end{proof}
\begin{definition}
    For $n\in\Z_{\geq0}$, we define,
\end{definition}
\begin{remark}
  Let $\cat$ be a $\K$-linear category and $M$ be a set of objects. Then the full-subcategory of $\cat$ of $M$ is generally neither additive nor idempotent complete.
  However, let $\langle M\rangle_{\oplus}$ be the full subcategory of $\cat$ on direct sums of objects from $M$. Then it is additive. Also $\langle M\rangle_{\stackrel{\subseteq}{\oplus}}$ the full subcategory of $\cat$ that 
  consists of on direct summands of direct sums of objects in $M$. It's additive and idempotent complete.
\end{remark}

\begin{remark}
    Let $\cat$ be an additive, idempotent complete $\K$-linear category. Let $M$ be a set of objects in $\cat$. Generally the full subcategory of $M$ is neither additive nor idempotent complete.
\end{remark}

\begin{definition}
    Let $\cat$ be an additive, idempotent, $\K$-linear category and $M$ be a set of objects in $\cat$. We define,
    \begin{enumerate}[1).]
        \item $\langle M\rangle_{\oplus}$ as the full subcategory of $\cat$ with objects direct sums of objects from $M$.
        \item $\langle M\rangle_{\oplus,\summand}$ as the full subcategory of $\cat$ with objects direct summands of direct sums of objects in $M$.
    \end{enumerate}
\end{definition}

\begin{proposition}
    $\langle M\rangle_\oplus$ is additive and $\langle M\rangle_{\oplus,\summand}$ is additive and idempotent complete.
\end{proposition}

\begin{proof}
    Additivity is clear in both cases. Let $X\in\langle M\rangle_{\oplus,\summand}$ and $e\in\End(X)$ be an idempotent. In $\cat$, we have 
    \begin{equation*}
        \varphi:X\stackrel{\cong}{\to} X_1\oplus X_2.
    \end{equation*}
    By definition, we can assume that $X_1,X_2\in\langle M\rangle_{\oplus,\summand}$. Thus it is idempotent complete.
\end{proof}

\begin{lemma}
  Let $\cat$ be a $\K$-linear category, we have the following statements.
  \begin{enumerate}[1).]
    \item $\cat^{\oplus},\cat^{\oplus,\summand}$ are additive $\K$-linear.
    \item $\cat^{\summand},\cat^{\oplus,\summand}$ are idempotent complete.
    \item $\cat$ is additive then $\cat\cong \cat^{\oplus}$ as $\K$-linear category.
    \item $\cat$ is idempotent complete then $\cat\cong \cat^{\summand}$ as $\K$-linear categories.
  \end{enumerate}
\end{lemma}

\begin{proof}
    For \textbf{1)}, the additiity is clear. For \textbf{2)}, let $(X,e)\in\cat^{\summand}$ and $efe\in\End((X,e))$ be an idempotent. Then $(X,efe)\in\cat^{\summand}$. Let
    \begin{equation*}
        p=efe\id_Xe\in\cat^{\summand}((X,e),(X,efe)),q=e\id_Xefe\in\cat^{\summand}((X,efe),(X,e)).
    \end{equation*}
    Then we have,
    \begin{equation*}
        pq = efeeeefe = efe \id_X efe = \id_{(X,efe)},qp = eefeefee = eefee=efe.
    \end{equation*}
    By Lemma \ref{lemma_idempotent_equivalent_condition}, this completes the proof.
    \par\textbf{3)}. is clear by setting $F:\cat\to\cat^{\oplus}$ be such that $F((X,p_i,q_i)) = (X_i)$. Remains to show that $\cat\cong\cat^{\summand}$ for idempotent complete category $\cat$. 
    Let $F:\cat\to\cat^{\summand}$ be such that $F(X)\to (X,\id_X)$. Clearly, this is fully-faithful. Remains to show this is essentially surjective. For $(X,e)\in\cat^{\summand}$, again by Lemma \ref{lemma_idempotent_equivalent_condition}, there is $Y\in\cat$ such that 
    \begin{equation*}
        \exists p:X\to Y,q:Y\to X, qp = e,pq = \id_Y.
    \end{equation*}
    Thus we have $pe\in\cat^{\summand}((X,e),(Y,\id_Y))$ and $eq\in\cat^{\summand}((Y,\id_Y),(X,e))$. Then 
    \begin{align*}
        peeq &= peq = pqpq = \id_Y=\id_{(Y,\id_Y)},\\
        eqpe & = e = \id_{(X,e)}.
    \end{align*}
\end{proof}

\end{document}